\scanpage{115}
% L1–3
Soit $\deleted{X}A$ un anneau commutatif, $X$ le schéma affine qu'il définit,
\[
 A^{\unsure{\mathrm{cons}|\mathrm{con}}}
 =\underset{\unsure{\mathscr{P}|P|\mathfrak P}}{\underline{\lim}}\ A_{\mathscr{P}}
\]
% L4–6
$\mathscr{P}$ parcourt les ``partitions \unsure{standard|standart}'' de $X$ en ens constructibles loc. fermés \unsure{,|et} affines, et pour une telle $\mathscr{P}=(X_i)_{i\in I}$, $A_{\mathscr{P}}$ désigne \deleted{\unsure{le produit}} $\unsure{A|\Gamma}(\coprod X_{i\,\mathrm{red}})=\deleted{\unsure{\mathscr{P}}}\prod A(X_{i\,\mathrm{red}})$

% L7–8
Les hom de transition sont injectifs. Considérons
\[
 \widetilde X=\operatorname{Spec}(A^{\mathrm{cons}})
\]
% L9
On a $\operatorname{Top}(\widetilde X)=\underset{\mathscr P}{\underleftarrow{\operatorname{Lim}}}\operatorname{Top}\operatorname{Spec}(A_{\mathscr P})$

% L10–11
En particulier
\[
 \operatorname{Ens}(\widetilde X)=\underset{\mathscr P}{\underleftarrow{\operatorname{Lim}}}\ 
 \underbrace{\operatorname{Ens}\operatorname{Spec} A_{\mathscr P}}_{\coprod X_i}
\]
% L12–14
et comme les morphismes de transition sont des \deleted{\unsure{homéomorphismes}}\inserted{bijections}, on trouve \deleted{que} une \deleted{\unsure{surj}}\inserted{bijection}.
\[
 \operatorname{Ens}\widetilde X\simeq\operatorname{Ens} X
\]
% L15–18
Cette bijection est compatible avec les topologies de $\widetilde X$ et celle de $X^{\mathrm{cons}}$, comme --- on voit aisément sur la description d'une base d'ouverts de $\widetilde X$ \ldots
% archive sheet number: 114
